\documentclass[12pt]{article}

\setlength{\parindent}{0pt}
\setcounter{secnumdepth}{3}
\usepackage[margin=1in]{geometry}

% packages
\usepackage{amsmath, amsfonts, amsthm, amssymb}
\usepackage{graphicx, float}
\usepackage{mathtools}
\usepackage{titlesec}
\usepackage{interval}
\usepackage{hyperref}
\usepackage{siunitx}
\usepackage{titling}
\usepackage{vwcol}
\usepackage{setspace}
\usepackage{empheq}
\usepackage{cancel}
\usepackage{esdiff}
\usepackage{multicol}
\usepackage{mdframed}
\usepackage{esdiff}
\usepackage{tikzsymbols}
\usepackage{multicol}
\usepackage{tikz}
\usepackage{varwidth}
\usepackage{pgfplots}
\usepackage{fancyhdr}
\usepackage{enumitem}
\usepackage{tcolorbox}

\intervalconfig {
	soft open fences
}

% header
\pagestyle{fancy}
\fancyhf{}
\lhead{Ethan Fan}
\rhead{AP Calculus BC}
\cfoot{\thepage}

% section format
\titleformat{\section}
  {\large \bfseries}
  {}
  {0em}
  {}
\titlespacing*{\section}{0pt}{0.5em}{0.3em}

\titleformat{\subsection}[runin]
  {\bfseries}
  {}
  {0em}
  {}
\titlespacing*{\subsection}{0pt}{0pt}{0em}

\titleformat{\subsubsection}[runin]
  {\itshape}
  {(\roman{subsubsection})}
  {0em}
  {}

% commands
\newcommand{\alignedintertext}[1]{%
  \noalign{%
    \vskip\belowdisplayshortskip
    \vtop{\hsize=\linewidth#1\par
    \expandafter}%
    \expandafter\prevdepth\the\prevdepth
  }%
}

\newcommand*{\paren}[1]{\ensuremath\left(#1\right)}
\newcommand*{\limit}[2][x]{\ensuremath{\displaystyle\lim_{#1\to#2}}}
\newcommand*{\Limit}[3][x]{\ensuremath{\displaystyle\lim_{#1\to#2}\left[#3\right]}}
\newcommand*{\deriv}[1][x]{\ensuremath{\dfrac{\mathrm{d}}{\mathrm{d}#1}}}
\newcommand*{\Deriv}[2][x]{\ensuremath{\dfrac{\mathrm{d}}{\mathrm{d}#1}\left[#2\right]}}
\newcommand{\ddx}{\dfrac{d}{dx}}
\newcommand{\dydx}{\dfrac{dy}{dx}}
\newcommand{\dydt}{\dfrac{dy}{dt}}
\newcommand{\dxdt}{\dfrac{dx}{dt}}
\newcommand*{\iinteg}[2][x]{\ensuremath{\displaystyle\int #2\;\mathrm{d}#1}}
\newcommand*{\dinteg}[4][x]{\ensuremath{\displaystyle\int_{#2}^{#3}#4\;\mathrm{d}#1}}
\newcommand*{\abs}[1]{\ensuremath{\left|#1\right|}}
\newcommand*{\eps}{\varepsilon}
\newcommand*{\real}{\ensuremath\mathbb{R}}
\newcommand*{\integer}{\ensuremath\mathbb{Z}}
\newcommand*{\posint}{\ensuremath\mathbb{Z^+}}
\newcommand*{\cbrt}[1]{\ensuremath{\sqrt[3]{#1}}}
\newcommand*{\lheqzero}{\ensuremath{\underset{\text{L'H}}{\overset{\left[\frac00\right]}{=}}}}
\newcommand*{\lheqinfty}{\ensuremath{\underset{\text{L'H}}{\overset{\left[\frac{\infty}{\infty}\right]}{=}}}}
\newcommand{\tor}{\text{ or }}
\newcommand{\floor}[1]{\lfloor #1 \rfloor}

\newcommand{\vem}{\vspace{0.5em}}
\newcommand{\example}[1]{\section{Example #1}}
\newcommand{\problem}[1]{\section{Problem #1}}
\newcommand{\subproblem}[1]{\subsection{(#1)} \,}

\DeclareMathOperator{\dne}{DNE}
\DeclareMathOperator{\sgn}{sgn}

\newcommand{\definition}[1]{\begin{tcolorbox}[colback=red!5!white,colframe=red!75!black,parbox=false] #1 \end{tcolorbox}}

\newcommand{\theorem}[2]{\begin{tcolorbox}[title={#1},colback=blue!5!white,colframe=blue!75!black,parbox=false] #2 \end{tcolorbox}}

\allowdisplaybreaks
% \postdisplaypenalty=100000

% document
\begin{document}

\vspace*{0cm}

% opening
\begin{center}
    {\Large \bfseries Problem Set 17} \\[0.5cm]
    {\large Partial Fractions} \\[0.35cm]
    {\normalsize October 16, 2026}
\end{center}

% problems

\problem{1}

\subproblem{b}
\begin{gather*}
    \int \frac{6x-2}{(x-3)(x+1)}\,dx =\int \frac{4}{x-3} + \frac{2}{x+1}\, dx\\
    = \boxed{4 \ln |x-3| + 2\ln |x+1|+C}
\end{gather*}

\subproblem{c}
\begin{gather*}
    \int \frac{2x+5}{(x+4)(x-2)} \, dx = \int \frac{3}{2(x-2)} + \frac{1}{2(x+4)}  \, dx\\
    = \boxed{ \frac{3}{2} \ln |x-2| + \frac{1}{2} \ln |x+4|+C}
\end{gather*}

\problem{2}

\subproblem{b}
\begin{gather*}
    \int \frac{x^2+4x-1}{x(x+1)(x-1)} \, dx = \int \frac{1}{x} + \frac{2}{x-1} - \frac{2}{x+1} \, dx\\
    = \boxed{\ln |x| + 2 \ln |x-1| -2 \ln|x+1|+C}
\end{gather*}

\subproblem{c}
\begin{gather*}
    \int \frac{4x+28}{(x+1)(x-1)(x-3)} \, dx = \int \frac{3}{x+1}- \frac{8}{x-1} +\frac{5}{x-3}\, dx\\
    = \boxed{ 3 \ln |x+1| - 8 \ln |x-1| + 5\ln |x-3|+C}
\end{gather*}

\problem{3}

\subproblem{b}
\begin{gather*}
    \int \frac{3x^2-2x-3}{x^2(x-1)}\, dx= \int \frac{5}{x} - \frac{2}{x-1} + \frac{3}{x^2}\, dx\\
    = \boxed{ 5 \ln |x| - 2\ln |x-1| - \frac{3}{x}+C}
\end{gather*}

\subproblem{c}
\begin{gather*}
    \int \frac{x^2}{(x+1)^3} \, dx = \int \frac{1}{x+1} - \frac{2}{(x+1)^2} + \frac{1}{(x+1)^3} \, dx\\
    = \boxed{\ln |x+1|+ \frac{2}{x+1}   - \frac{1}{2(x+1)^2} +C} 
\end{gather*}

\subproblem{d}
The suggested form fails to work, as it can be simplified into a single factor. Three constants must exist, each corresponding to one of the three degrees of the denominator's polynomial.

\problem{4}

\subproblem{b}
\begin{gather*}
    \int \frac{x^2+x+1}{(x+1)(x^2+1)}  \, dx= \int \frac{1}{2(x+1)} + \frac{x+1}{2(x^2+1)}\, dx\\
    = \boxed{ \frac{1}{2} \ln |x+1| + \frac{1}{4} \ln |x^2+1| + \frac{1}{2}\arctan x+C}
\end{gather*}

\problem{5}
\begin{gather*}
    \int \ln (x^2-x+2)\, dx = \begin{bmatrix}
        u = \ln (x^2-x+2) & dv =dx\\
        du = \frac{2x-1}{x^2-x+2} \, dx& v=x
    \end{bmatrix}\\ 
    x\ln (x^2-x+2) - \int \frac{2x^2-x}{x^2-x+2}\, dx\\
    \frac{2x^2-x}{x^2-x+2}= 2+\frac{x-4}{x^2-x+2} = 2+ \frac{\frac{1}{2}(2x-1) - \frac{7}{2}}{x^2-x+2}\\
    \int \ln (x^2-x+2)\, dx = x\ln (x^2-x+2) - 2x \\- \frac{1}{2} \ln (x^2-x+2) + \frac{7}{2} \int \frac{1}{(x-\frac{1}{2})^2 + \frac{7}{4}} \, dx\\
    = \boxed{\paren{x-\frac{1}{2}} \ln (x^2-x+2) -2x + \sqrt{7 } \arctan \paren{\frac{2x-1}{\sqrt{7}}}+C}
\end{gather*}
$x^2-x+2$ is not reducible.

\problem{6}
\begin{gather*}
    \int_1^2 \frac{x^2+1}{3x-x^2}\, dx = \int_1^2 -1+ \frac{3x+1}{x(3-x)} \, dx\\
    =-1 + \int_1^2 \frac{1}{3x} + \frac{10}{3(3-x)}\, dx\\
    =-1 + (\frac{1}{3} \ln x- \frac{10}{3} \ln (3-x)) \Big|^2_1\\
    = \boxed{-1  + \frac{11}{3} \ln 2}
\end{gather*}

\problem{7}

\subproblem{a}
\begin{gather*}
    \sin \frac{\theta}{2}=\frac{t}{\sqrt{1+t^2}},\, \cos \frac{\theta}{2}= \frac{1}{\sqrt{1+t^2}}\\
    \sin \theta = 2 \cdot \frac{t}{\sqrt{1+t^2}} \cdot \frac{1}{\sqrt{1+t^2}} = \boxed{\frac{2t}{1+t^2}}\\
    \cos \theta = \paren{\frac{1}{\sqrt{1+t^2}}}^2 - \paren{\frac{t}{\sqrt{1+t^2}}}^2 = \boxed{\frac{1-t^2}{1+t^2}}\\
    \theta = 2 \arctan t \implies \boxed{d \theta = \frac{2\, dt}{1+t^2}}
\end{gather*}

\subproblem{b}
\begin{gather*}
    \int_0^\frac{\pi}{2} \frac{dx}{1+\sin x } =  \int_0^1 \frac{2\, dt}{t^2+2t+1}\\
    = -\frac{2}{(t+1)} \Big|^1_0 = 2-1=\boxed{1}
\end{gather*}

\subproblem{c}
\begin{gather*}
    \int \frac{dx}{3\sin x -4\cos x} = \int \frac{1}{2t^2+3t-2} \, dt\\
    =\int \frac{2}{5(2t-1)} - \frac{1}{5(t+2)} \, dt\\
    = \boxed{ \frac{1}{5} \ln |2t-1| - \frac{1}{5} \ln |t+2|+C}
\end{gather*}

\problem{8}

\subproblem{a}
\begin{align*}
    T &= \int_{100}^{900} \frac{dP}{0.0002 P(1000-P)} \\
    &= 5000\int_{100}^{900} \frac{1}{1000P} +\frac{1}{1000(1000-P)} \,dP\\
    &= 5 \ln |P| - 5\ln |1000-P| \Big|^{900}_{100}\\
    &= \boxed{20 \ln 3}
\end{align*}

\subproblem{b}
\begin{gather*}
    T_1=T_2=10\ln 3
\end{gather*}
The two integrals have the same value as the function is symmetric about $P=500$.

\problem{9}

\subproblem{a}
\begin{gather*}
    \sum_{k=1}^{99} \frac{1}{k(k+1)}= \sum_{k=1}^{99} \frac{1}{k}- \frac{1}{k+1} = \frac{1}{1}-\frac{1}{100} = \boxed{\frac{99}{100}}
\end{gather*}

\subproblem{b}
\begin{gather*}
    \sum_{k=1}^{n} \frac{1}{k(k+2)} =  \frac{1}{2}\sum_{k=1}^{n} \frac{1}{k} - \frac{1}{k+2} \\
    =\frac{1}{2} \paren{1+\frac{1}{2}- \frac{1}{n+1}-\frac{1}{n+2}}\\
    = \boxed{\frac{3}{4}- \frac{2n+3}{2(n+1)(n+2)}}\\
    \lim_{n \to \infty} \sum_{k=1}^{n} \frac{1}{k(k+2)} = \boxed{\frac{3}{4}}
\end{gather*}

\problem{10}

\subproblem{a}
\begin{gather*}
    \frac{1}{Q(x)} = \sum_{i=1}^n \frac{A_i}{x-a_i}  \iff 1= \sum_{i=1}^n A_i\prod_{j \ne i} (x-a_j)\\
    \sum_{i=1}^n A_i \prod_{j \ne i} (x-a_j) \Big|_{x=a_k} = A_k \prod_{j \ne k} (a_k-a_j) \Big|_{x=a_k}=\frac{A_k}{A_k} = 1
\end{gather*}

\subproblem{b}
\begin{gather*}
    \frac{x}{Q(x)} =  \sum_{i=1}^n \frac{A_ix}{x-a_i}\\
    \lim_{x\to \infty} \frac{x}{Q(x)}=0\\
    \lim_{x\to \infty}\frac{A_ix}{x-a_i} = A_i\lim_{x\to \infty}\frac{1}{1-\frac{a_i}{x}} = A_i = 0\\
    \boxed{\sum_{i=1}^n A_i = 0}
\end{gather*}

\subproblem{c}
\begin{gather*}
    \frac{x^k}{Q(x)} =  \sum_{i=1}^n \frac{a_i^k}{(x-a_i) \prod_{j \ne i} (a_i-a_j) }\\
    \frac{x^{k+1}}{Q(x)} =  \sum_{i=1}^n \frac{a_i^k}{ \prod_{j \ne i} (a_i-a_j) }  \cdot \frac{x}{(x-a_i)}\\
    0 \le k \le n-2 \implies \lim_{x \to \infty} \frac{x^{k+1}}{Q(x)} = 0\\
    \iff \lim_{x \to \infty} \sum_{i=1}^n \frac{a_i^k}{ \prod_{j \ne i} (a_i-a_j) }  \cdot \frac{x}{(x-a_i)} = 0\\
    k=n-1 \implies \lim_{x \to \infty} \frac{x^{k+1}}{Q(x)} = 1\\
    \iff \lim_{x \to \infty} \sum_{i=1}^n \frac{a_i^k}{ \prod_{j \ne i} (a_i-a_j) }  \cdot \frac{x}{(x-a_i)} = 1
\end{gather*}

\subproblem{d}
\begin{gather*}
    \sum_{i=1}^n \frac{a_i^k}{ \prod_{j \ne i} (a_i-a_j) } = \sum_{i=1}^n \frac{1}{ \prod_{j \ne i} (i-j) }=0 \\
    \prod_{j \ne i} (i-j) = (i-1)!(n-i)!(-1)^{n-i}\\
    \sum_{i=1}^n \frac{(-1)^{n-i}}{ (i-1)!(n-i)! } = 0\\
    \iff \sum_{m=0}^{n-1} \frac{(-1)^{n-m-1}}{ m!(n-m-1)! } = 0\\
    \iff \sum_{m=0}^{n-1} (-1)^{m}\frac{(n-1)!}{ m!(n-m-1)! } = 0\\
    \iff \boxed{\sum_{m=0}^{n-1} (-1)^{m} \binom{n - 1}{m} = 0}
\end{gather*}

\subproblem{e}
\begin{gather*}
    Q(x)=x(x-1)(x-2)(x-3)\\
    a_1=0:A_1= -\frac{1}{6}\\
    a_2=1:A_2= \frac{1}{2}\\
    a_3=2: A_3=-\frac{1}{2}\\
    a_4=3: A_4= \frac{1}{6}\\
    \sum_{i=1}^4 A_i = -\frac{1}{6} + \frac{1}{2} - \frac{1}{2} + \frac{1}{6} = 0\\
    \int \frac{dx}{x(x-1)(x-2)(x-3)}\\
     = \int \left( \frac{-1/6}{x} + \frac{1/2}{x-1} - \frac{1/2}{x-2} + \frac{1/6}{x-3} \right) dx \\
     = \boxed{-\frac{1}{6} \ln\vert{}x\vert{} + \frac{1}{2} \ln\vert{}x-1\vert{} - \frac{1}{2} \ln\vert{}x-2\vert{} + \frac{1}{6} \ln\vert{}x-3\vert{} + C}
\end{gather*}







































\end{document}

